\section{Discussion}
\label{sec:discussion}

\subsection{Why visits and conversions disagree}
\label{sec:why-disagree}

The same learners, data, and protocol give opposite verdicts on the two
outcomes. Two mechanisms plausibly combine, and we cannot cleanly separate
them with this dataset.

\begin{enumerate}
  \item \textbf{Signal.} Conversions are 0.29\,\% of users. The conversion
        effect scores are far less stable across seeds (rank correlation
        0.49 versus 0.98), so an uplift ranking for conversions is mostly
        estimation noise layered on a signal that a simple response model
        already captures.
  \item \textbf{Structure.} Independently of noise, the share of conversions
        that would have happened anyway is flat (about 67\,\%) across score
        groups. When lift scales with baseline propensity, the response
        ranking and the uplift ranking are nearly the same ranking, and
        there is nothing for uplift to add.
\end{enumerate}

For visits the picture reverses: the highest-probability visitors are mostly
``sure things'' (88\,\% would have visited anyway), and a model that sees
treatment can move budget to users the ad actually changes. This matches the
dataset authors' advice to model uplift on visits.

\subsection{Randomization checks are not optional}
\label{sec:imbalance}

The pooled benchmark is advertised as randomized, and the headline
difference in means would be reported as such. A classifier two-sample test
and a joint likelihood-ratio test show that treatment is weakly but
detectably predictable, and that the imbalance correlates with the outcomes.
Correcting it moves the visit effect by 28\,\%. The observational
demonstration shows the adjusted estimators do recover the answer when a
known confounder is injected, which supports, but does not prove, that the
adjustment is doing the right thing on the real imbalance.

\subsection{Assignment, exposure, and what ``effect'' means}
\label{sec:assignment}

The intent-to-treat effect (assignment) and the effect of seeing an ad differ
by more than an order of magnitude on visits, because 96.4\,\% of eligible
users are never shown an ad. A targeting policy acts on eligibility, so the
ATE and uplift results above are the operationally relevant ones; the
instrumental-variable estimate answers a different question, about users who
do get exposed.

\subsection{Limitations}
\label{sec:limitations}

\paragraph{The data is a benchmark, not a live system.} Non-uniform
sub-sampling, including negative sampling on labels, means absolute rates and
effects are not real-world ad incrementality, and the business values in the
targeting stage are assumptions.

\paragraph{Adjustment has limits.} Treatment is only conditionally ignorable
here, and adjusted estimators assume the 12 features capture the dependence
between assignment and outcomes. If the sub-sampling depended on outcomes
themselves, no feature adjustment fully corrects it. The mechanism behind
the imbalance (per-test pooling or label-dependent sampling) is a
hypothesis, not verified.

\paragraph{Anonymized features.} Features are randomly projected, so
heterogeneity findings describe feature segments, not interpretable user
types.

\paragraph{Exposure estimates rest on the exclusion restriction.} The IV
estimates assume assignment affects outcomes only through ads shown.

\paragraph{Subsampled models.} The random forest, logistic regression, and
causal forest were trained on 1--2M-row samples for time. LightGBM and the
meta-learners used all 8.39M training rows. The causal forest's weaker
conversion result may partly reflect this.

\paragraph{One dataset, one split.} Every result comes from a single
stratified split. Seed and disjoint-half refits address estimator variance
but not variation across alternative splits or other advertisers' data.
